\section{Regression}
\label{sec:regression}

\subsection{The regression problem}

In a \emph{regression} problem we want to predict a real-valued quantity from
some input. The input is described by a \emph{feature vector}
$x \in \X \subseteq \R^d$, and the quantity we want to predict is a
\emph{response} $y \in \R$.

\begin{example}
Suppose we want to predict the price of a flat. Each flat can be encoded by a
feature vector
\[
x = \big( [\text{surface in m}^2],\ [\text{number of rooms}],\ [\text{floor}],\ \dots \big) \in \R^d ,
\]
and the response $y \in \R$ is its price. Each coordinate of $x$ is a
\emph{feature} (also called \emph{covariate} or \emph{explanatory variable}),
and $d$ is the number of features.
\end{example}

We are given a \emph{training set} of $n$ pairs
\[
(x_1, y_1), \dots, (x_n, y_n) \in \R^d \times \R ,
\]
where $n$ is the number of training samples and $d$ the dimension of each data
point. As in the rest of these notes, we stack the inputs into the
\emph{design matrix} and the outputs into a single vector:
\[
X = \begin{pmatrix}
- & x_1 & - \\
  & \vdots &  \\
- & x_n & -
\end{pmatrix} \in \R^{n \times d},
\qquad
y = (y_1, \dots, y_n) \in \R^n .
\]

\begin{remark}[Encoding non-numerical features]
We have implicitly assumed that each input is a vector of real numbers,
$x \in \R^d$. In practice many features are \emph{categorical} rather than
continuous: a person's sex (\texttt{male}/\texttt{female}), a colour, a city,
etc. We turn them into numbers as follows.
\begin{itemize}
    \item A \emph{binary} feature can be encoded by a single $-1/1$ indicator, e.g.\
    $\texttt{male} \mapsto -1$ and $\texttt{female} \mapsto 1$.
    \item A categorical feature with $K \geq 3$ possible values should
    \emph{not} simply be mapped to $1, 2, \dots, K$: this would invent a
    spurious order and notion of distance between categories (it would suggest
    that ``city $3$'' lies between ``city $2$'' and ``city $4$''). Instead one
    uses \emph{one-hot} (or \emph{dummy}) encoding: the feature is replaced by
    $K$ binary features, exactly one of which equals $1$. For example the colour
    $\{\texttt{red}, \texttt{green}, \texttt{blue}\}$ becomes
    \[
    \texttt{red} \mapsto (1,0,0), \quad
    \texttt{green} \mapsto (0,1,0), \quad
    \texttt{blue} \mapsto (0,0,1) .
    \]
    One might object that for $K = 2$ the one-hot encoding does not give back the "$-1/1$"
    encoding of the previous item. It in fact does; as explained later in the notes.
\end{itemize}
After such an encoding every input is again a genuine vector $x \in \R^d$, and
the whole framework below applies unchanged.
\end{remark}

The goal is to \emph{learn} a function $f : \X \to \R$ such that
\[
f(x) \approx y \qquad \text{for unseen pairs } (x, y).
\]
Once such an $f$ is found, we can predict the response of a brand new input
$x$ simply by evaluating $f(x)$.

\paragraph{What do we need to specify?}
Turning the vague wish ``$f(x) \approx y$'' into a well-posed problem requires
two ingredients:
\begin{enumerate}
    \item a \emph{class of functions} $\F$ in which we look for $f$ (e.g.\ all
    linear functions of $x$);
    \item a way to \emph{quantify} how good a candidate $f$ is on the data,
    i.e.\ a \emph{loss function}.
\end{enumerate}

\subsection{Choosing the function class}
\label{subsec:function-class}

The function class $\F$ collects all the candidate predictors we are willing to
consider. In practice we never search over an abstract set of functions:
instead we fix a \emph{parametrisation}, that is, we describe every candidate by
a finite vector of \emph{parameters} (or \emph{weights}) $w \in\R^p$ (where $p$ corresponds to the number of trainable parameters) and denote
the corresponding predictor by $f_w$. The function class is then
\[
\F = \big\{\, f_w \ : \ w \in \R^p \,\big\},
\]
and \emph{learning the function reduces to searching for a good parameter}
$w$.  How $w$ enters $f_w$ is a modelling choice; here are the main examples. (For
the linear models that occupy the rest of these notes the parameter has exactly one
coordinate per feature, so it lives in the same $\R^d$ as the inputs; the last two
examples below are parametrised differently, and there $w$ simply collects however
many weights the model has.)

\paragraph{Linear in the input.}
The simplest choice makes the prediction a linear function of the input,
\[
f_w(x) = \langle w, x \rangle = \sum_{j=1}^d w_j x_j , \qquad w \in \R^d ,
\]
so that the parameter $w$ has the same dimension $d$ as the input $x$ (i.e., $p = d$).


\paragraph{Nonlinear in the parameters.}
Nothing forces $w$ to enter linearly. One may equally consider parametrisations
in which the parameter appears \emph{inside} a nonlinear function: the resulting objective is non-convex in $w$, and must be tackled by
iterative optimisation (e.g.\ gradient descent), with no guarantee of reaching a
global minimum. The most prominent examples are neural networks.

\paragraph{Neural networks.}
A prominent and very expressive parametrisation that is nonlinear in $w$ is a
\emph{neural network}. In its simplest, one-hidden-layer form,
\[
f_w(x) = \sum_{k=1}^m a_k \, \sigma\big( \langle b_k, x\rangle + c_k \big),
\]
where $\sigma$ is a fixed nonlinearity (e.g.\ a ReLU or a sigmoid) and the
parameters $w = (a_k, b_k, c_k)_{k=1}^m$ gather all the weights. Increasing the
width $m$ (or stacking several such layers) yields an extremely rich
function class. The price to pay is, once more, a highly non-convex dependence
on $w$, which is the central difficulty of modern deep learning. 


\subsection{Choosing the loss function}

Fixing a function class is not enough: we still need to say what it means for a
prediction to be \emph{good}. This is the role of a \emph{loss function}
$\ell(\hat{y}, y)$, which quantifies how much we penalise predicting
$\hat{y} = f(x)$ when the true response is $y$. Averaging the loss over the
training set gives the \emph{empirical risk}, which we view as a function of the
parameter $w$ and denote
\[
L(w) = \frac{1}{n} \sum_{i=1}^n \ell\big( f_w(x_i), y_i \big).
\]
\emph{Empirical risk minimisation} then consists in choosing the parameter $w$
(equivalently, $f_w \in \F$) so as to make $L(w)$ as small as possible.

\begin{remark}[The constant $1/n$ does not matter for the minimiser]
Multiplying an objective by a positive constant does not move its minimisers:
for any $c > 0$, $\argmin_w c\, L(w) = \argmin_w L(w)$.
In particular the factor $1/n$ in front of the sum is immaterial when our only
goal is to \emph{locate} the best $w$. We nonetheless keep it throughout these
notes, because it makes $L(w)$ an \emph{average} loss per sample rather than a
total: its scale then does not grow with the size of the training set, and losses
computed on datasets of different sizes remain comparable. Be aware, though, that
the constant \emph{does} rescale the gradient $\nabla L$, and hence matters once a
step size enters the picture, as in gradient descent below. (The extra $\tfrac12$
that will appear for the squared error is a separate matter: it belongs to the
loss $\ell$ itself and is there only to cancel the $2$ coming from the derivative.)
\end{remark}

\paragraph{Several ways to penalise an error.}
The choice of $\ell$ is itself a modelling decision: different losses encode
different ideas of what a good fit is. A few common choices are
\begin{itemize}
    \item the \emph{squared error} $\ell(\hat{y}, y) = \tfrac{1}{2}(\hat{y} - y)^2$,
    by far the most common, which strongly penalises large residuals;
    \item the \emph{absolute error} $\ell(\hat{y}, y) = |\hat{y} - y|$, more
    robust to outliers but non-differentiable at~$0$;
    \item the \emph{Huber loss}, a compromise that is quadratic for small
    residuals and linear for large ones.
\end{itemize}
In the remainder of this section we focus on the squared error, which, when
combined with a model that is linear in $w$, yields the particularly clean
\emph{least-squares} problem.

\begin{remark}[Training error is not the goal]
Making the empirical risk small only guarantees that $f$ fits the
\emph{training} data. What we actually care about is the \emph{test error},
i.e.\ how well $f$ predicts on new, unseen inputs. A function can fit the
training set perfectly and still generalise poorly.
\end{remark}


\subsection{Minimising the empirical risk}

Choosing a function class and a loss has turned learning into an optimisation
problem:
\[
w^\star \in \argmin_{w \in \R^p} \ L(w),
\]
where $\argmin_w L(w)$ denotes the set of parameters at which $L$
attains its minimum, which is not necessarily unique (and may not even exist as we will see later on in the classification setting).

For most models and losses there is no formula for the minimiser, so we
approximate $w^\star$ by an \emph{iterative} method that only needs to evaluate
gradients of $L$. The basic idea is \emph{gradient descent}: from some starting
point $w_0$, repeatedly take a small step opposite to the gradient,
\[
w_{t+1} = w_t - \gamma\, \nabla L(w_t),
\]
and its cheaper \emph{stochastic} variant, which at each step uses the gradient at
only a few sampled data points. We will come back to both in detail in
\Cref{subsec:gd}; for now it is enough to keep in mind that minimising $L$
is generally done by such gradient-based iteration.


\section{Least squares}
\label{sec:least-squares}

We now specialise the two ingredients above to the simplest case:
\[
f_w(x) = \langle w, x \rangle = \sum_{j=1}^d w_j x_j , \qquad w \in \R^d ,
\]
together with the squared-error loss. The function class is thus
\[
\F = \big\{\, x \mapsto \langle w, x \rangle \ : \ w \in \R^d \,\big\} .
\]

\begin{remark}[Adding an intercept]
Often we want an \emph{affine} predictor $f(x) = \langle w, x \rangle + b$ with
an intercept (or bias) $b \in \R$, so that $f(0) \neq 0$ is allowed. This fits
into the same framework: simply append a constant coordinate equal to $1$ to
each input, $\tilde{x} = (x, 1) \in \R^{d+1}$, and a coordinate $b$ to the
weights, $\tilde{w} = (w, b) \in \R^{d+1}$. Then
$\langle \tilde{w}, \tilde{x}\rangle = \langle w, x \rangle + b$, and we are
back to a purely linear model in dimension $d+1$. For this reason we keep the
notation $f_w(x) = \langle w, x \rangle$ below without loss of generality.
\end{remark}

\paragraph{The least-squares objective.}
Plugging the linear model into the empirical risk with the square loss, learning amounts to
minimising over $w \in \R^d$ the function
\[
L(w) = \frac{1}{2n} \sum_{i=1}^n \big( \langle w, x_i \rangle - y_i \big)^2
     = \frac{1}{2n} \, \Vert X w - y \Vert^2 ,
\]
where the second equality follows from stacking the residuals
$\langle w, x_i\rangle - y_i$ into the vector $Xw - y \in \R^n$. 

\begin{remark}
Notice that $L$ is a positive function, and that $L(w) = 0$ if and only if $X w = y$, i.e. if and only if the vector $w$ perfectly fits the dataset: $\langle x_i, w \rangle = y_i$ for all $i \in \{1, \dots, n\}$.
\end{remark}

\noindent
The problem
\[
w^\star \in \argmin_{w \in \R^d} \ \frac{1}{2n}\, \Vert X w - y \Vert^2
\]
is called the \emph{least-squares} problem.

\begin{exercise}[$\star \star$]
Check that we indeed have that 
\[
\frac{1}{2n} \sum_{i=1}^n \big( \langle w, x_i \rangle - y_i \big)^2
   = \frac{1}{2n} \, \Vert X w - y \Vert^2 .
\]
Then write $L$ as an explicit quadratic in $w$:
\[
L(w) = \tfrac{1}{2n}\, w^\top X^\top X\, w - \tfrac{1}{n}\, w^\top X^\top y + \tfrac{1}{2n}\, \Vert y \Vert^2 .
\]
Deduce that $\nabla L(w) = \tfrac{1}{n}\, X^\top (X w - y)$ and that the Hessian is
$\nabla^2 L(w) = \tfrac{1}{n}\, X^\top X$, independently of $w$.
\end{exercise}

\paragraph{Solving the least squares problem.}
The function $L$ is convex (it is a quadratic with positive semi-definite
Hessian $\tfrac{1}{n} X^\top X \succeq 0$, see \Cref{sec:cheat-convexity}), so any
point where the gradient vanishes is a global minimum. Computing the gradient of $L$ (see
\Cref{sec:cheat-linalg}) leads to:
\[
\nabla L(w) = \tfrac{1}{n}\, X^\top (X w - y) .
\]
Setting $\nabla L(w^\star) = 0$, the factor $\tfrac{1}{n}$ cancels and we obtain the
\emph{normal equations}
\[
\ X^\top X \, w^\star = X^\top y  .
\]

\noindent
From here, we distinguish two different settings:
\begin{itemize}
    \item The \emph{under-parametrised} setting (sometimes called the \textit{over-determined} setting), where there are fewer trainable parameters than data samples, i.e.\ $d \leq n$. In this setting, we will see that the normal equations admit a \emph{unique} solution\footnote{Under the mild assumption that $X$ has full column rank.}, and hence $L$ has a unique global minimiser.
    \item The \emph{over-parametrised} setting (sometimes called the \textit{under-determined} setting), where there are more trainable parameters than data samples, i.e.\ $d > n$. In this setting, we will see that the normal equations admit \emph{infinitely many} solutions, and hence $L$ has infinitely many global minimisers.
\end{itemize}

\subsection{Underparametrised setting ($n \geq d$): a unique solution}
\label{subsec:underparam}

Solving the normal equations for $w^\star$ requires inverting $X^\top X \in
\R^{d\times d}$, so before stating the closed-form solution we need to
understand when this is possible.

\paragraph{In the underparametrised setting, $X^\top X$ is (nearly always) invertible.}
From \Cref{app:lemma:rank-XX-X}, we have that $X^\top X$ is invertible if and only if $X$ has full column rank, i.e. if and only if $\spn(x_1, \dots, x_n) = \R^d$. This means that the inputs must span the whole $\R^d$ space, hence $n \geq d$. In fact, one can show that for ``nearly all'' datasets for which $n \geq d$, it holds that $\rank(X) = d$ and hence $X^\top X$ is invertible (see \Cref{subsec:cheat-rank} for more details). 

\begin{prop}[Closed-form solution]\label{prop:closed-form-LS}
Assuming that $\spn(x_1, \dots, x_n) = \R^d$, then $X^\top X \in \R^{d \times d}$ is invertible and the least-squares problem
has a unique minimiser, given by
\[
w^\star = (X^\top X)^{-1} X^\top y .
\]
\end{prop}

\paragraph{Geometric interpretation.}
The vector of predictions $X w$ ranges, as $w$ varies over $\R^d$, exactly over
the column space $\spn(\text{columns of } X) \subseteq \R^n$. Minimising
$\Vert Xw - y\Vert$ therefore means finding the point of this subspace that is
\emph{closest} to $y$: the optimal prediction $X w^\star$ is the
\emph{orthogonal projection} of $y$ onto the column space of $X$.
Equivalently, the residual $X w^\star - y$ is orthogonal to every column of
$X$, which is precisely the statement of the normal equations
$X^\top (X w^\star - y) = 0$.


\begin{exercise}[$\star\star$]
Assume $X^\top X$ is invertible, show that $L$ can be rewritten around its global
minimiser $w^\star$ as:
\[
L(w) = L(w^\star) + \tfrac{1}{2n}\, \big\Vert X (w - w^\star) \big\Vert^2.
\]
\end{exercise}




\subsection{Overparametrised setting ($d > n$): an infinite space of solutions}
\label{subsec:overparam}

We now consider the overparametrised setting where $d > n$. In this case $X^\top X$ cannot be inverted, so we solve the normal equations differently. 

\paragraph{Generic interpolation.}
In the overparametrised setting where $d > n$\footnote{in fact, all of the following also holds for $d=n$}, the inputs $x_1, \dots, x_n$ can span a linear subspace of $\R^d$ of dimension at most $n$. In fact, if they are linearly independent (which holds for almost all datasets, see \Cref{subsec:cheat-rank}), then their span is of dimension exactly $n$, which means that $\rank(X) = n$. This observation leads to the following proposition.

\begin{prop}[Zero training error]
If the inputs $x_1, \dots, x_n$ are linearly independent (which requires $d \geq n$), then $\rank(X) = n$ and the linear map
\[
X : \R^d \to \R^n, \qquad w \mapsto Xw
\]
is surjective. Hence for \emph{every} response vector $y \in \R^n$ there exists at least one
weight $w^\star$ with $Xw^\star = y$. This means that $y_i = \langle x_i, w^\star \rangle$ for all $i \in [n]$, which means that the
training data is fitted perfectly. Also note that this means that $L(w^\star) = \tfrac{1}{2n}\Vert Xw^\star - y\Vert^2 = 0$, hence that $w^\star \in \arg\min L$.
\end{prop}






\paragraph{The minimisers form an affine space.}
As soon as one interpolating weight exists the set of global
minimisers of $L$ is \emph{exactly} the set of interpolants,
\[
S = \big\{\, w \in \R^d \ : \ Xw = y \,\big\} .
\]
This is an \emph{affine subspace} of $\R^d$. Now let $w^\star$ be any particular
solution ($X w^\star = y$), then any other solution $w$ such that $Xw = y$ is equivalent to $X (w - w^\star) = 0$, which means that $\langle x_i, w - w^\star \rangle = 0$ for every $i$. Hence $w - w^\star$ is orthogonal to every input, that is, $w - w^\star \in \spn(x_1, \dots, x_n)^{\perp}$. Hence 
\[
S = w^\star + \spn(x_1, \dots, x_n)^{\perp}.
\]
Now assuming that the inputs are linearly independent, then the dimension of the affine space $S$ is equal to $d-n$. In other words, there is a whole $(d-n)$-dimensional affine space of weights achieving
zero training error.


\paragraph{The minimum-norm solution.}
Because the whole affine space $S$ is optimal for the training loss, we need an
extra criterion to pick a single predictor. A natural and very common choice is
the \emph{minimum-norm} interpolant, i.e. the point of $S$ closest to the origin:
\[
w^\dagger = \argmin_{w \,:\, Xw = y} \ \Vert w \Vert .
\]

\begin{prop}[Minimum-norm solution]
If $\rank(X) = n$, then $XX^\top \in \R^{n \times n}$ is invertible and
\[
w^\dagger = X^\top (XX^\top)^{-1} y
\]
is the unique minimum-norm solution of $Xw = y$.
\end{prop}

\begin{proof}[Geometric proof]
The idea is purely geometric: $w^\dagger$ is the orthogonal projection of the
origin onto the affine space of interpolants, and Pythagoras does the rest. Recall
that the set of minimisers of $L$ is the affine space
\[
S = \{\, w \in \R^d : Xw = y \,\} ,
\]
whose direction is the subspace $\spn(x_1, \dots, x_n)^\perp$.

\emph{Step 1: $w^\dagger$ is an interpolant lying in $\spn(x_1, \dots, x_n)$.}
Since $\rank(X) = n$, the matrix $XX^\top$ is invertible, so $w^\dagger = X^\top
(XX^\top)^{-1} y$ is well defined and
\[
X w^\dagger = XX^\top (XX^\top)^{-1} y = y ,
\]
so $w^\dagger \in S$. Moreover, writing $u = (XX^\top)^{-1} y \in \R^n$, we have
$w^\dagger = X^\top u = \sum_{i=1}^n u_i\, x_i$, so $w^\dagger \in \spn(x_1, \dots,
x_n)$.

\emph{Step 2: Pythagoras.}
Let $w^\star \in S$ be an arbitrary interpolant, i.e.\
$X w^\star = y$. Decompose it around $w^\dagger$ as
\[
w^\star = w^\dagger + h, \qquad h := w^\star - w^\dagger .
\]
Because $X w^\star = X w^\dagger = y$, the difference satisfies $X h = 0$, that is,
$h \in \spn(x_1, \dots, x_n)^\perp$. But $w^\dagger \in \spn(x_1, \dots, x_n)$ by Step 1, so $w^\dagger$ and $h$
are orthogonal, $\langle w^\dagger, h \rangle = 0$. The Pythagorean theorem then
gives
\[
\Vert w^\star \Vert^2 = \Vert w^\dagger + h \Vert^2
= \Vert w^\dagger \Vert^2 + \Vert h \Vert^2 \geq \Vert w^\dagger \Vert^2 ,
\]
so $\Vert w^\star \Vert \geq \Vert w^\dagger \Vert$ for every interpolant $w^\star$.
Equality forces $\Vert h \Vert = 0$, i.e.\ $w^\star = w^\dagger$, which proves that
$w^\dagger$ is the \emph{unique} minimum-norm interpolant.
\end{proof}



\begin{remark}[Two dual regimes]
The formulas $w^\star = (X^\top X)^{-1} X^\top y$ (under-parametrised, $n \geq d$)
and $w^\dagger = X^\top (XX^\top)^{-1} y$ (over-parametrised, $d > n$) are mirror
images: in the first $X^\top X \in \R^{d \times d}$ is invertible and the solution
is unique; in the second $XX^\top \in \R^{n \times n}$ is invertible while $X^\top
X$ is singular, so the solution is only unique \emph{after} imposing the
minimum-norm rule. Both are special cases of the \emph{Moore-Penrose
pseudo-inverse} $X^+$, for which $w = X^+ y$ in either regime.
\end{remark}

\begin{exercise}[$\star$]
In the boundary regime where $n = d$ with $\rank(X) = n = d$, prove that the under-parametrised and over-parametrised formulas match, i.e. that
$(X^\top X)^{-1} X^\top y = X^\top (X X^\top)^{-1}y$.
\end{exercise}

% \begin{remark}[Interpolation is not the end of the story]
% Fitting the training data perfectly says nothing about the \emph{test} error. In
% the over-parametrised regime the choice of which interpolant we land on (here, the
% minimum-norm one) is exactly what controls generalisation, a theme we return to
% with the bias-variance trade-off.
% \end{remark}


\subsection{Gradient descent and stochastic gradient descent}
\label{subsec:gd}

The above equations give a closed-form solution to least squares, but computing
them is not always practical, for two reasons.
\begin{itemize}
    \item \textbf{Cost.} In practice one does not compute the inverse
    $(X^\top X)^{-1}$ (nor $(X X^\top)^{-1}$)  explicitly, which is both wasteful and numerically unstable.
    The normal equations $X^\top X\, w = X^\top y$ form a linear system, and one
    solves it directly, for instance by a Cholesky factorisation of the symmetric
    positive definite matrix $X^\top X$, or, better conditioned, by a $QR$ factorisation or
    an SVD of $X$ that avoids forming $X^\top X$ at all. The cost is nonetheless of
    order $O(n d^2 + d^3)$. When the number of features $d$ is large, of the order
    of thousands or millions as in modern models, this becomes prohibitive in both
    time and memory, whichever method is used.
    \item \textbf{Generality.} The closed form is special to the squared loss with
    a model that is linear in $w$. As soon as we change the loss (absolute error,
    Huber, \dots) or use a model that is nonlinear in $w$ (a neural network), there
    is \emph{no} closed-form minimiser at all.
\end{itemize}
In both cases we give up on solving for $w$ exactly and instead compute an
approximate minimiser \emph{iteratively}, using only gradients of $L$.

\paragraph{Gradient descent.}
Gradient descent (GD) starts from an arbitrary $w_0 \in \R^d$ and repeatedly takes
a small step in the direction of steepest descent, i.e.\ opposite to the gradient:
\[
w_{t+1} = w_t - \gamma\, \nabla L(w_t), \qquad t = 0, 1, 2, \dots,
\]
where $\gamma > 0$ is the \emph{step size} (or \emph{learning rate}). The
intuition is that $-\nabla L(w_t)$ points in the direction along which $L$
decreases the most, so for a small enough $\gamma$ each step lowers the loss. For
the least-squares objective the gradient was computed above,
\[
\nabla L(w) = \tfrac{1}{n}\, X^\top (X w - y),
\]
so one GD step costs a matrix-vector product with $X$ and with $X^\top$, that is
$O(n d)$ operations, and requires no matrix inversion. Since $L$ is convex, GD with
a suitable step size converges to a global minimiser.

\begin{remark}[Choosing the step size]
\label{rem:step-size}
The step size $\gamma$ must be tuned. If $\gamma$ is too large the iterates can
overshoot the minimum and even diverge; if it is too small the algorithm is
stable but crawls, needing very many iterations. For least squares one can be
precise, because the Hessian $\nabla^2 L = \tfrac{1}{n} X^\top X$ is a constant
matrix. Write
\[
\beta \coloneqq \lambda_{\max}\big(\tfrac{1}{n} X^\top X\big) = \tfrac{1}{n}\lambda_{\max}(X^\top X)
\]
for its largest eigenvalue, the \emph{smoothness constant} of $L$. Then the iterates
converge for every step size in the range
\[
0 < \gamma < \frac{2}{\beta} ,
\]
and diverge as soon as $\gamma > 2/\beta$;
this is exactly the threshold appearing in
\Cref{fact:gd-converges} below. The usual default is $\gamma = 1/\beta$, which sits
comfortably inside the admissible range.
\end{remark}

\paragraph{A variational viewpoint on gradient descent.}
Gradient descent looks like a purely mechanical rule, ``take a step against the
gradient''. It is enlightening to notice that each step is in fact the solution of a
small optimisation problem: at $w_t$ we do not minimise $L$ itself (which we cannot
do in closed form), but a cheap \emph{surrogate} obtained by linearising $L$ around
$w_t$ and forbidding ourselves from moving too far.

\begin{prop}[GD as linearise-and-penalise]
\label{prop:gd-variational}
For any $w_t \in \R^d$ and step size $\gamma > 0$, the gradient descent update
$w_{t+1} = w_t - \gamma\, \nabla L(w_t)$ is the unique minimiser of the surrogate
objective
\[
w_{t+1} = \argmin_{w \in \R^d} \ \Big\{
\underbrace{L(w_t) + \langle \nabla L(w_t),\, w - w_t \rangle}_{\text{linearisation of } L \text{ at } w_t}
\ + \ \frac{1}{2\gamma} \Vert w - w_t \Vert^2 \Big\} .
\]
\end{prop}

\begin{proof}
The bracketed function is a strictly convex quadratic in $w$ (the linear part
contributes nothing to the curvature, and the penalty $\tfrac{1}{2\gamma}\Vert w - w_t
\Vert^2$ has positive-definite Hessian $\tfrac{1}{\gamma} I$), so it has a unique
minimiser, found by cancelling its gradient:
\[
\nabla L(w_t) + \frac{1}{\gamma}\, (w - w_t) = 0
\quad \Longleftrightarrow \quad
w = w_t - \gamma\, \nabla L(w_t) . \qedhere
\]
\end{proof}

So the step size $\gamma$ plays the role of a \emph{trust region}: a large $\gamma$
weakens the proximity penalty and lets the step follow the linear model further,
while a small $\gamma$ keeps $w_{t+1}$ close to $w_t$, where the linearisation is
trustworthy.

\begin{exercise}[$\star$]
What recursion do we obtain if we keep the true loss $L$ instead of its
linearisation, i.e.\ if we define $w_{t+1} = \argmin_{w} \big\{ L(w) +
\tfrac{1}{2\gamma}\Vert w - w_t \Vert^2 \big\}$? (This is known as \emph{implicit}
gradient descent, or the proximal point method.)
\end{exercise}

\paragraph{Stochastic gradient descent.}
Even though GD avoids matrix inversion, every single step still evaluates the full
gradient, which is a sum over all $n$ training points. Writing the empirical risk
as a finite sum,
\[
L(w) = \frac{1}{n} \sum_{i=1}^n \ell_i(w), \qquad
\ell_i(w) = \ell\big( f_w(x_i), y_i \big),
\]
we see that one full gradient $\nabla L(w) = \frac1n \sum_{i=1}^n \nabla \ell_i(w)$
costs $O(n d)$, which is again expensive when the number of samples $n$ is huge.

\emph{Stochastic gradient descent} (SGD) replaces the full gradient by the
gradient at a \emph{single} training point, drawn at random. At each step one
samples an index $i_t$ uniformly in $\{1, \dots, n\}$ and updates
\[
w_{t+1} = w_t - \gamma_t\, \nabla \ell_{i_t}(w_t).
\]
The key point is that this cheap direction is, on average, the right one:
\[
\E_{i_t}\big[ \nabla \ell_{i_t}(w_t) \big]
   = \frac{1}{n} \sum_{i=1}^n \nabla \ell_i(w_t) = \nabla L(w_t),
\]
so $\nabla \ell_{i_t}(w_t)$ is an \emph{unbiased estimate} of the true gradient,
obtained at cost $O(d)$ instead of $O(n d)$. For least squares the per-sample
gradient is simply
\[
\nabla \ell_i(w) = \big( \langle w, x_i \rangle - y_i \big)\, x_i .
\]

\paragraph{The vocabulary you will meet in practice.}
Two conventions surround SGD in every library, and they are worth naming.

\emph{Mini-batches.} Rather than a single sample, one usually averages the gradient
over a small \emph{mini-batch} $B_t \subseteq \{1, \dots, n\}$ of $B = |B_t|$ points
drawn at random,
\[
w_{t+1} = w_t - \gamma_t \, \frac{1}{B} \sum_{i \in B_t} \nabla \ell_i(w_t) .
\]
This is still an unbiased estimate of $\nabla L(w_t)$, costs $O(Bd)$, and its variance
is smaller by a factor $B$. The two extremes $B = 1$ and $B = n$ are plain SGD and
plain GD, so $B$ interpolates between them, trading variance against cost per step.
Typical values are $B = 32$ to $512$: large enough for the step to be a meaningful
direction, small enough that many steps are taken per pass through the data.

\emph{Epochs.} One \emph{epoch} is one full pass over the training set, that is
$n / B$ steps. Training time is reported in epochs rather than in steps, because it is
the number of times each data point has been looked at. In practice the mini-batches
of an epoch are usually obtained by shuffling the $n$ indices once and cutting the
result into blocks of size $B$ (sampling \emph{without} replacement), which is
marginally biased compared with the idealised uniform sampling above, but simpler and
slightly more efficient.


\begin{exercise}[$\star$]
What is the analogue of \Cref{prop:gd-variational} for SGD? That is, which
surrogate objective does the SGD update $w_{t+1} = w_t - \gamma\, \nabla
\ell_{i_t}(w_t)$ minimise at each step?
\end{exercise}

\paragraph{Implicit bias of gradient descent.}
In the over-parametrised regime $d > n$ there is a whole affine space of
minimisers of $L$, and it is a priori unclear which one will be recovered by gradient descent. To answer this question, the crucial observation is that \emph{every gradient of $L$ lives in
the span of the data}. Indeed,
\[
\nabla L(w) = \tfrac{1}{n} X^\top (X w - y) = \tfrac{1}{n} \sum_{i=1}^n
\big( \langle w, x_i \rangle - y_i \big)\, x_i \ \in \ \spn(x_1, \dots, x_n).
\]
Write $V \coloneqq
\spn(x_1, \dots, x_n)$ for this subspace.

Throughout, we take for granted the following statement. It is a genuine
\emph{theorem}, not a hypothesis on the data: it can be but it lies outside the scope of this class, so we simply admit it here.

\begin{fact}[Convergence of the iterates, admitted]
\label{fact:gd-converges}
For the least-squares loss and a small enough step size (any $0 < \gamma < 2/\beta$
with $\beta = \tfrac1n \lambda_{\max}(X^\top X)$ the largest eigenvalue of the
Hessian), the gradient descent iterates converge to some limit $w_\infty =
\lim_{t \to \infty} w_t$.
\end{fact}

\begin{prop}[Implicit bias of GD]
Let $(w_t)_{t \geq 0}$ be the gradient descent iterates $w_{t+1} = w_t - \gamma\,
\nabla L(w_t)$ on the least-squares loss $L$, started at $w_0 \in \R^d$, and let
$w_\infty$ be their limit. Then:
\begin{enumerate}
    \item $w_\infty$ is a minimiser of $L$;
    \item among \emph{all} minimisers of $L$, the point $w_\infty$ is the one closest
    to the starting point $w_0$, that is, the orthogonal projection of $w_0$ onto the
    set of solutions:
    \[
    w_\infty = \argmin_{w \in S}\ \Vert w - w_0 \Vert .
    \]
    In particular, starting from $w_0 = 0$ recovers the minimum-norm solution
    $w^\dagger$.
\end{enumerate}
\end{prop}

\begin{proof}[Geometric proof]
\emph{(1) The iterates stay in $w_0 + V$.} Each gradient $\nabla L(w_t)$ lies in
$V$ by the observation above, so $w_{t+1} - w_t = -\gamma\, \nabla L(w_t) \in V$. By
induction $w_t - w_0 \in V$ for every $t$, i.e.\ $w_t \in w_0 + V$. Since $V$ is a
closed subspace, the limit inherits the property: $w_\infty \in w_0 + V$.

\emph{(2) $w_\infty$ is a minimiser.} This is the only place the convergence
assumption is used. Passing to the limit in the update $w_{t+1} = w_t - \gamma\,
\nabla L(w_t)$ and using continuity of $\nabla L$, the limit is a fixed point,
$w_\infty = w_\infty - \gamma\, \nabla L(w_\infty)$, which forces $\nabla L(w_\infty)
= 0$. As $L$ is convex, this means $w_\infty \in \argmin L =: S$, the
set of solutions of the normal equations. Recall that $S$ is an affine space with
direction $V^\perp$.

\emph{(3) Pythagoras.} Let $w^\star \in S$ be any minimiser. Both $w_\infty$ and
$w^\star$ lie in $S$, whose direction is $V^\perp$, so $w_\infty - w^\star
\in V^\perp$. On the other hand $w_\infty - w_0 \in V$ by part (1). These two vectors
are therefore orthogonal, and Pythagoras gives
\[
\Vert w^\star - w_0 \Vert^2
= \Vert \underbrace{(w^\star - w_\infty)}_{\in\, V^\perp}
      + \underbrace{(w_\infty - w_0)}_{\in\, V} \Vert^2
= \Vert w^\star - w_\infty \Vert^2 + \Vert w_\infty - w_0 \Vert^2
\ \geq\ \Vert w_\infty - w_0 \Vert^2 .
\]
Thus $w_\infty$ minimises the distance to $w_0$ over all minimisers, with equality
only when $w^\star = w_\infty$: it is the orthogonal projection of $w_0$ onto $S$.
\end{proof}

This is what is meant by the \emph{implicit bias} (or \emph{implicit
regularisation}) of gradient descent: although the loss $L$ is blind to the choice
of interpolant, the optimisation algorithm is not, and it silently selects the
solution closest to its initialisation.

\begin{exercise}[$\star$]
What happens for SGD? And what if the loss $\ell$ is not the square norm, but some
other coercive loss?
\end{exercise}


\subsection{Ridge regression (weight decay)}
\label{subsec:ridge}

The implicit bias of gradient descent selects a solution \emph{implicitly}, through
the dynamics of the algorithm. Ridge regression does the same thing
\emph{explicitly}, by adding a penalty to the loss itself that discourages large
weights. One augments the least-squares objective with an $\ell_2$ penalty (called
\emph{weight decay} in the deep-learning community),
\[
L_\lambda(w) = \underbrace{\frac{1}{2n}\Vert X w - y \Vert^2}_{L(w)}
\ + \ \frac{\lambda}{2}\Vert w \Vert^2 ,
\]
where $\lambda > 0$ is the \emph{regularisation strength}. The extra term pulls the
solution towards the origin; the larger $\lambda$, the stronger the pull.

\begin{remark}[Two practical warnings: the intercept and the scale of the features]
Unlike everything that came before, the penalty $\Vert w \Vert^2$ is \emph{not}
invariant under the harmless-looking reparametrisations we have been allowing
ourselves, and this has two concrete consequences.
\begin{itemize}
    \item \emph{Do not penalise the intercept.} Recall that we absorbed an intercept
    $b$ into $w$ by appending a constant coordinate $1$ to each input. If we now
    penalise the whole of $\Vert w \Vert^2$, we are also shrinking $b$ towards $0$,
    which amounts to saying that a response centred around $1000$ is somehow less
    plausible than one centred around $0$: the fit would change if we merely chose to
    measure $y$ in different units, or from a different origin. In practice one
    therefore penalises only the genuine weights, minimising
    $\tfrac{1}{2n}\Vert Xw + b\mathbf{1} - y\Vert^2 + \tfrac{\lambda}{2}\Vert w\Vert^2$,
    or equivalently centres $y$ and the columns of $X$ beforehand and drops the
    intercept altogether.
    \item \emph{Standardise the features first.} A single penalty $\lambda$ is applied
    to every coordinate of $w$, but the coordinates of $x$ need not be comparable: a
    surface in $\text{m}^2$ and a surface in $\text{km}^2$ differ by a factor $10^6$,
    and the corresponding weight by a factor $10^{-6}$. Ridge would then penalise the
    two encodings of the \emph{same} model completely differently, whereas ordinary
    least squares is entirely insensitive to such a rescaling. The remedy, and the
    other reasons for applying it, are collected in \Cref{subsec:standardise}.
\end{itemize}
Below we ignore both points and penalise $\Vert w \Vert^2$ as written, which keeps
the formulas clean; just remember to centre and standardise before applying them.
\end{remark}

Unlike plain least squares, ridge regression
always has a \emph{unique} solution, even in the over-parametrised regime $d > n$
where $X^\top X$ is not invertible.

\begin{prop}[Existence, uniqueness, closed form]
For every $\lambda > 0$ the objective $L_\lambda$ is strictly convex and coercive, so
it admits a \emph{unique} minimiser
\[
w^\star_\lambda = \big( X^\top X + n\lambda I \big)^{-1} X^\top y ,
\]
which is well defined for \emph{any} $X \in \R^{n \times d}$, whatever the relative
sizes of $n$ and $d$.
\end{prop}

\begin{proof}
The Hessian of $L_\lambda$ is $\nabla^2 L_\lambda = \tfrac{1}{n} X^\top X + \lambda I
\succeq \lambda I \succ 0$, so $L_\lambda$ is strictly convex (see
\Cref{sec:cheat-convexity}); and $L_\lambda(w) \geq \tfrac{\lambda}{2}\Vert w
\Vert^2 \to \infty$ as $\Vert w \Vert \to \infty$, so it is coercive and its infimum
is attained. A strictly convex function has at most one minimiser, hence exactly one
here. Cancelling the gradient,
\[
\nabla L_\lambda(w) = \tfrac{1}{n} X^\top (X w - y) + \lambda w = 0
\quad \Longleftrightarrow \quad
\big( X^\top X + n\lambda I \big)\, w = X^\top y .
\]
Finally $X^\top X + n\lambda I$ is positive definite (its eigenvalues are $\geq
n\lambda > 0$), hence invertible, which gives the announced formula.
\end{proof}

It is worth noting that $w^\star_\lambda$ always lies in the span of the data: from
the optimality condition, $w^\star_\lambda = \tfrac{1}{n\lambda} X^\top (y - X
w^\star_\lambda) \in \spn(x_1, \dots, x_n)$.

\paragraph{The two extremes.}
The parameter $\lambda$ interpolates between two meaningful limits: crank it up and
the penalty dominates, killing the weights; send it to zero and we recover ordinary
least squares, more precisely its minimum-norm solution.

\begin{prop}[Limits in $\lambda$]
Let $w^\dagger$ denote the minimum-norm minimiser of $L$ (the orthogonal projection
of $0$ onto the affine set $S = \argmin L$). Then
\[
\lim_{\lambda \to +\infty} w^\star_\lambda = 0
\qquad \text{and} \qquad
\lim_{\lambda \to 0^+} w^\star_\lambda = w^\dagger .
\]
\end{prop}

\begin{proof}
\emph{Limit $\lambda \to \infty$.} Comparing the optimal value with the one at $w =
0$,
\[
\tfrac{\lambda}{2}\Vert w^\star_\lambda \Vert^2 \leq L_\lambda(w^\star_\lambda)
\leq L_\lambda(0) = L(0) = \tfrac{1}{2n}\Vert y \Vert^2 ,
\]
so $\Vert w^\star_\lambda \Vert^2 \leq \Vert y \Vert^2 / (n\lambda) \to 0$.

\emph{Limit $\lambda \to 0^+$.} By optimality of $w^\star_\lambda$ for $L_\lambda$,
tested against $w^\dagger$,
\[
L(w^\star_\lambda) + \tfrac{\lambda}{2}\Vert w^\star_\lambda \Vert^2
\ \leq\ L(w^\dagger) + \tfrac{\lambda}{2}\Vert w^\dagger \Vert^2 .
\tag{$\ast$}
\]
Since $w^\dagger$ minimises $L$, we have $L(w^\star_\lambda) \geq L(w^\dagger)$, and
($\ast$) gives $\Vert w^\star_\lambda \Vert \leq \Vert w^\dagger \Vert$: the family is
bounded. Rearranging ($\ast$) also gives $0 \leq L(w^\star_\lambda) - L(w^\dagger)
\leq \tfrac{\lambda}{2}\Vert w^\dagger \Vert^2 \to 0$, so $L(w^\star_\lambda) \to \min
L$. Now let $\bar w$ be the limit of any convergent subsequence $w^\star_{\lambda_k}$
(one exists, by boundedness). By continuity $L(\bar w) = \min L$, so $\bar w \in S$,
and $\Vert \bar w \Vert \leq \Vert w^\dagger \Vert$. But $w^\dagger$ is the
\emph{unique} element of $S$ of minimal norm, so $\bar w = w^\dagger$. Every
convergent subsequence thus has the same limit $w^\dagger$, and since the family is
bounded, $w^\star_\lambda \to w^\dagger$.
\end{proof}

\begin{remark}[Two roads to the same solution]
The limit $\lambda \to 0^+$ is remarkable: it recovers exactly the minimum-norm solution $w^\dagger$ that gradient descent from $w_0 = 0$ finds through its implicit
bias. Explicit regularisation (ridge, taken to the zero-penalty limit) and implicit
regularisation (gradient descent) therefore lead to the \emph{same} predictor. This
is one of the cleanest illustrations of why $w^\dagger$ is the ``natural'' least-squares
solution in the over-parametrised regime.
\end{remark}

\subsection{Feature expansion}
\label{subsec:features}

\paragraph{A linear model can be a curved function.}
Models which are linear in the input data $x$ are often too rigid. A flexible remedy is to transform the input using a \emph{chosen} (as opposed to learnt) transformation $\phi$. Such a function is often called a \emph{feature map}:
\[
\phi : \R^{d} \to \R^{d'}, \qquad x \mapsto \phi(x) = \big( \phi_1(x), \dots, \phi_{d'}(x) \big),
\]
which sends the raw input, living in its native dimension $d$, to a vector of
$d'$ new features. After this \emph{feature augmentation}, one wants to learn a linear model of these features:
\[
f_w(x) = \langle w, \phi(x) \rangle,
\qquad w \in \R^{d'} .
\]
The number of features $d'$ is ours to choose, and it may be larger, sometimes
much larger, than the raw input dimension $d$.

\begin{example}[Polynomial features]
\label{ex:poly-features}
For a scalar input $z \in \R$ (so $d = 1$), the feature map
$\phi(z) = (1, z, z^2, \dots, z^k) \in \R^{k+1}$ turns the model into a
degree-$k$ polynomial
\[  
f_w(z) = w_0 + w_1 z + \dots + w_k z^k ,
\]
built out of $d' = k+1$ features. The predictor is now a \emph{nonlinear}
function of $z$, yet it remains a \emph{linear} function of the parameter $w$.
\end{example}




\paragraph{Everything happens as if the data had been transformed once and for all.}
Transform each input once, before doing anything else, and call the result
$\tilde{x}_i = \phi(x_i) \in \R^{d'}$. We are then fitting the plain linear
model $\tilde{x} \mapsto \langle w, \tilde{x} \rangle$ on the dataset
$(\tilde{x}_1, y_1), \dots, (\tilde{x}_n, y_n)$, whose design matrix stacks the
\emph{transformed} inputs:
\[
\tilde{X} = \begin{pmatrix}
- & \tilde{x}_1 & - \\
  & \vdots &  \\
- & \tilde{x}_n & -
\end{pmatrix} = \begin{pmatrix}
- & \phi(x_1) & - \\
  & \vdots &  \\
- & \phi(x_n) & -
\end{pmatrix} \in \R^{n \times d'},
\qquad
y = (y_1, \dots, y_n) \in \R^n .
\]
The empirical risk is then
\[
L(w) = \frac{1}{2n} \sum_{i=1}^n \big( \langle w, \phi(x_i) \rangle - y_i \big)^2
 = \frac{1}{2n} \sum_{i=1}^n \big( \langle w, \tilde{x}_i \rangle - y_i \big)^2
     = \frac{1}{2n} \big\Vert \tilde{X} w - y \big\Vert^2 ,
\]
which is word for word the least-squares problem of \Cref{sec:least-squares},
with $\tilde{X}$ in place of $X$ and $d'$ in place of $d$. Normal equations,
closed-form solution, minimum-norm interpolant, gradient descent and its
implicit bias, ridge regression: not one of these formulas ever cares whether
the rows of the design matrix were transformed beforehand. 


\paragraph{Under and over-parametrisation with polynomial expansion.}
Feature augmentation lets us increase $d'$ freely, and polynomial features are
the cleanest way to see what happens as $d'$ crosses $n$, the number of
training points. For a scalar input $z \in \R$, write $\phi^{(k)}$ for the
feature map of degree $k$,
\[
\phi^{(0)}(z) = (1), \quad
\phi^{(1)}(z) = \begin{pmatrix} 1 \\ z \end{pmatrix}, \quad
\phi^{(2)}(z) = \begin{pmatrix} 1 \\ z \\ z^2 \end{pmatrix}, \quad \dots, \quad
\phi^{(k)}(z) = \begin{pmatrix} 1 \\ z \\ \vdots \\ z^k \end{pmatrix} \in \R^{k+1} .
\]
A degree-$k$ polynomial model therefore works in
dimension $d' = k+1$, and nothing stops us from taking $d'$ as large as we like. Given $n$
training points $(z_1, y_1), \dots, (z_n, y_n)$ with distinct abscissae $z_i \in \R$, three regimes appear:
\begin{itemize}
    \item \emph{(Strictly) under-parametrised setting $d' < n$ (polynomials of degree less than $n - 1$)}: it is highly unlikely to be able to exactly pass through $n$ points with a polynomial of degree less than $n - 1$. There exists a unique polynomial which minimises the square error.
    \item \emph{Boundary $d' = n$ (polynomials of degree $n - 1$)}: there exists a unique polynomial of degree $(n-1)$ passing through all of the points exactly (Lagrange interpolation). That is, a weight vector $w \in \R^n$ with $f_w(x_i) = y_i$ for every $i$.
    \item \emph{Over-parametrised setting $d' > n$ (polynomials of degree $\geq n$)}: there exists an infinite number of ways of interpolating the points.
\end{itemize}

\begin{remark}[What ``linear regression'' really means]
``Linear regression'' refers to models that are \emph{linear in the parameters}
$w$, \emph{not} necessarily linear in the input $x$. Both $f_w(x) = \langle
w, x\rangle$ and the polynomial model $f_w(x) = \langle w, \phi(x)\rangle$ above
are linear regression, even though the latter is a curved function of $x$. This
is exactly the property that makes the least-squares objective a convex
quadratic in $w$, and it is why everything proved so far survives untouched, as
we now check.
\end{remark}

\paragraph{Learnt features: back to neural networks.}
Feature expansion also gives an interesting way of seeing the neural network we met in
\Cref{subsec:function-class}. Grouping its hidden units into a vector, one can write:
\[
f_w(x) = \sum_{k=1}^m a_k \, \sigma\big( \langle b_k, x \rangle + c_k \big)
= \big\langle a, \phi_{b,c}(x) \big\rangle ,
\qquad
\phi_{b,c}(x) = \Big( \sigma\big( \langle b_k, x \rangle + c_k \big) \Big)_{1 \leq k \leq m} ,
\]
we see that a neural network is nothing but a linear model on top of a feature
map: the $m$ hidden neurons are the features, and the output weights $a \in
\R^m$ play the role of $w$. If the inner parameters $(b_k, c_k)$ are frozen, for
instance drawn at random once and never touched again, then $\phi_{b,c}$ is a
fixed feature map and fitting $a$ is exactly the least-squares problem of this
section, closed-form solution included. This is the \emph{random features}
model.

The whole difference is that a neural network does not freeze them: it trains
the $(b_k, c_k)$ together with $a$. The features are then no longer picked by
hand beforehand, they are \emph{learnt from the data}, which is what makes these
models so effective on inputs (images, sound, text) for which nobody knows how
to design a good $\phi$. The price is the one announced earlier: the model is no
longer linear in its parameters, the objective is no longer convex, and there is
no closed-form solution.